% ------------------------------------------------------------------------------
% §1 Introduction — drafted last, Phase C
% ------------------------------------------------------------------------------
\section{Introduction}\label{sec:intro}

When an electricity market cannot keep the lights on by itself, the system
operator reaches outside it: a \emph{direction} orders a generator to run for
system security, and a compensation rule decides what that order pays.
Between 2022 and 2024, South Australia's direction scheme paid three gas
steam units at Torrens Island \$141.9 million for directed output the energy
market valued at \$7.2 million --- twenty dollars through the direction
channel for every dollar of market value. This paper asks whether that
transfer was an unfortunate price for security, or whether the scheme's own
design induced the scarcity it paid to relieve. The evidence supports the
second reading: the units' capacity withholding responds to the size of the
administered compensation price itself --- a margin to which no ordinary
market-power story assigns any role.

The setting has three ingredients that together make it a clean laboratory
for studying how out-of-market payments shape in-market conduct. First, a
\emph{security requirement}: South Australia's renewables-dominated grid
needs a minimum set of synchronous machines online, and when market dispatch
leaves the set incomplete, the operator directs specific units online. Second, a
\emph{gross payment}: a directed unit is compensated for its entire directed
output at a regulated reference price --- measured in this paper's data,
compensation is $0.95 \times$ directed MWh $\times$ the reference price, with
$R^{2}=0.99$ --- regardless of what it would have earned otherwise. Third, a
\emph{backward-looking price}: the reference price is the 90th percentile of
the trailing year of spot prices, so the 2022 energy crisis mechanically
paid itself forward into 2023--2024, holding the payment near \$350/MWh while
fuel costs halved. A unit that is loss-making in the energy market --- spot
exceeded fuel cost in 3.9\% of directed intervals --- faces a stark
choice between committing capacity to a market that will not cover its costs
and being absent, essential, directed, and paid gross at last year's scarcity
price.

The paper's empirical object is the availability margin of conduct: the
share of registered capacity a unit offers at prices at which it could
actually run. Two questions structure the analysis. RQ1: do units withhold
more when they are \emph{essential} --- when no combination of rival units
satisfies the security requirement without them? RQ2, the identifying
question: conditional on essentiality, does withholding respond to the
\emph{size} of the compensation price? The distinction carries the
identification. Being essential bundles compensation-eligibility with
whatever ordinary market power comes with being needed, and both predict
withholding-when-essential. But ordinary market power is indifferent to the
compensation formula; payment-seeking is not. A dose response to the
administered price, specific to essential intervals, is a signature only the
payment channel produces.

The answers. On RQ1, the directed workhorse station withholds more when
essential ($-0.076$ of registered capacity, stable in sign across all 72
leave-one-month-out re-estimates), and a direct control for energy-market
competition absorbs essentially none of the response --- indeed the
market-power alternative fails on the sign of its own conditions, because
essential periods are measurably \emph{more} competitive than ordinary ones,
not tighter. On RQ2, the essential-versus-matched withholding gap widens by
5.1 percentage points of registered capacity per \$100/MWh of compensation
price (wild-cluster-bootstrap $p<0.01$ on both outcome definitions), an
estimate whose interpretation was committed before it was run, and stable
across every treatment of the June 2022 crisis month. Decomposing the
outcome sharpens the finding onto the margin the mechanism actually pays:
the entire response is the probability that the unit's minimum-stable
quantum is within dispatch's reach --- the self-commitment choice that
averts or invites a direction --- which falls by 8.6 percentage points per
\$100/MWh (randomization-inference $p=0.046$), while offered depth
conditional on that choice responds not at all. The causal reading is
disciplined by the formula's own timing: gas prices fell nine months before
the payment did, and the gap tracked the payment, not the fuel, through
that wedge --- though the sharpest such contrast rests on four months and
clears only the 10\% level, so this paper's claim for the cross-month
gradient is ``consistent with payment-seeking,'' stated exactly.

Mechanism evidence locates the conduct precisely. The absence is a
\emph{standing posture}, not a daily calibration: the units' day-ahead bids
withhold or price out their floor megawatts on 76--80\% of all cleanly
observed days, in multi-day blocks, organised as a rotating station roster
that never once --- in 73 opportunities --- declared less than the security
standard required. Day-level tests, their interpretations committed in
advance and run on bids formed
outside direction exposure, find the posture indifferent to the daily
security envelope (evening exits are event-timed, not need-timed) and to the
daily prize (the offered floor price does not move on essential days). What
the direction buys is presence: directed output equals the units' operating
floor, with 55.9\% of episodes delivering zero or negative excess energy
over it. And the
paper reports its own correction: an earlier day-grain finding of
state-dependent withdrawal ($-12.7$ points on essential days, $p<0.001$) was
a contamination artifact --- 54\% of essential days had bids formed while a
direction was running or pending --- and vanishes on clean days. What
survives decontamination is the standing posture, its gross payment, and the
cross-month dose response.

The paper contributes to three literatures. The empirical market-power
literature in electricity measures conduct on the pricing margin: markups
over marginal cost \citep{Wolfram1999_duopoly,
BorensteinBushnellWolak2002_california}, bidding in multi-unit auctions
against theoretical benchmarks \citep{vonderFehrHarbord1993_spot,
HortacsuPuller2008_ercot}, pivotal-supplier ability and incentive
\citep{Wolak2003_unilateral, McRaeWolak2014_unilateral}, and physical
withholding --- as an input to price manipulation
\citep{JoskowKahn2002_california} or disguised as plant failures
\citep{FogelbergLazarczyk2019_failures,
BerglerHeimHuschelrath2017_failures}. The conduct documented here operates
on a margin that the literature treats as exogenous --- whether capacity is
offered at all --- and responds to a price that appears in none of its
models: the administered compensation rate. The nearest analytical
neighbour is the inc-dec gaming literature, in which units position in the
energy market in anticipation of out-of-market redispatch payments
\citep{HirthSchlecht2020_incdec}; this paper supplies field evidence of
that mechanism's signature, in a setting where the payment rate moves
through a mechanical formula. Second, the out-of-market payments
literature: \citet{BushnellWolak1999_rmr} argued that the design of
California's reliability-must-run contracts created serious incentive
problems and distorted bidding, and \citet{TeirilaRitz2019_capacity}
document market power in an administered capacity auction; this paper
supplies the payment-gradient evidence that debate never had. Third, the
capacity-mechanism design literature
\citep{CramtonOckenfelsStoft2013_capacity, Fabra2018_capacity}: the episode
is a natural experiment in what happens when an availability payment is
grossed up, backward-indexed, and layered onto an energy-only market ---
the design pays for absence, and absence is what it gets. Against all
three, one innocent account deserves naming early: start-up costs and
unit-commitment frictions rationalise block-level absence for steam units
\citep{Reguant2014_complex}, and this paper's mechanism evidence is fully
consistent with commitment costs shaping the posture's \emph{form}. What
commitment costs cannot produce is the posture's \emph{price sensitivity}
--- there is no unit-commitment reason for the self-commitment margin to
track a compensation formula --- which is exactly the object the paper's
identifying tests isolate.

The findings' scope is stated as carefully as their content. Nothing here
establishes daily intent: in this system the weather clock and the direction
clock are the same clock, and a desk watching either looks identical. The
claim is about the mechanism, not the desk's state of mind --- under this
payment structure a standing absent posture dominates commitment across
essentially all conditions, so rational conduct converges on exactly what the
data show, whatever its motive. That is an indictment of the design, and the
design is the policy lever: Section~\ref{sec:discussion} maps each defect
(gross payment, backward-looking price, unpriced presence) to a standard
repair, and notes the epilogue --- in September 2025 the synchronous
requirement was cut to one unit, closing the regime this paper studies.

Section~\ref{sec:setting} details the institutions; Section~\ref{sec:data}
the data and their validation against the operator's published statistics;
Section~\ref{sec:descriptives} the descriptive anatomy of the scheme;
Section~\ref{sec:design} the outcome, the essentiality flag, and the two
specifications; Section~\ref{sec:results} the regression results;
Section~\ref{sec:mechanism} the mechanism evidence and the contamination
correction; Section~\ref{sec:robustness} robustness;
Section~\ref{sec:discussion} concludes with design implications.
